\section{Local forms and the orbit oscillation estimate}
\label{sec:transfer}

We prove Proposition~\ref{prop:orbit-oscillation} from the local estimate
for inputs depending on spatial depth.  The proof constructs admissible
kernels for the difference of two smooth scales, averages over translations
of one dyadic lattice, and then transfers the resulting estimate first to
integer sequences and then to finite orbit segments.

\subsection{Admissible local forms}

Dyadic lattice intervals are taken half-open.  If \(J\) is a lattice
interval and \(I\subseteq J\) is a descendant, its depth in \(J\) is
\(s(I)=\log_2(|J|/|I|)\).  A finite rooted tree contains every ancestor
between each of its intervals and its root.

The kernel class is the one used in \cite[Section~2]{H}.
\begin{definition}[Admissible local kernels]
\label{def:transfer-local-kernels}
Fix \(c_0>0\) and \(C_0\geq1\).  For an interval \(I\) of length
\(r=|I|\), a kernel \(w_I\in C_c^\infty(\mathbb R^2;\mathbb C)\) is
admissible if the following hold.
\begin{enumerate}
\item For every \((x,t)\in\operatorname{supp}w_I\) and
 \(j\in\{0,1,2,3\}\),
 \(x+jt\in I\) and
 \(\operatorname{dist}(x+jt,\partial I)\geq c_0r\).
\item For every \(v\in\mathbb R\) and \(j=0,1,2,3\),
 \[
  \int_{\mathbb R}w_I(v-jt,t)\,dt=0.
 \]
\item For all nonnegative integers \(p,q\), with \(n=p+q\),
 \[
  \|\partial_x^p\partial_t^q w_I\|_\infty
     \leq r^{-n}\bigl(C_0(n+2)\bigr)^{C_0(n+2)}.
 \]
\end{enumerate}
For measurable scalar functions \(z_j:\mathbb R\to\mathbb C\), define
the associated local form, whenever its integral is finite, by
\[
H_I(z_0,z_1,z_2,z_3)
   =r^{-2}\int_{\mathbb R^2}w_I(x,t)
                  \prod_{j=0}^3z_j(x+jt)\,dx\,dt.
\]
For a Hilbert-valued argument with a distinguished physical scalar
coordinate, the form uses that coordinate.  We write
\(\|z\|_{p,I}=(|I|^{-1}\int_I\|z\|^p)^{1/p}\) for the normalized interval
norm, with the usual essential supremum when \(p=\infty\).  Norms without
an interval subscript use ordinary Lebesgue measure.
\end{definition}

For a nonempty consecutive integer block \(B=\{a,\ldots,b\}\) and a
sequence of complex functions \(h_s\) on any common domain, write
\[
 \mathcal V_B(h)(y)
   =\max_{s\in B}|h_s(y)|
       +\sum_{s=a}^{b-1}|h_{s+1}(y)-h_s(y)|.
\]
The sum is zero for a singleton block.

\begin{theorem}[Local estimate with variation in depth]
\label{thm:depth-local}
Fix the admissibility parameters \(c_0,C_0\).  Let \(J\) be an interval
in any dyadic lattice, let \(\mathcal T\) be a finite rooted tree below
\(J\), and let \(d\) be its largest depth in \(J\).  Let
\(\mathcal I\subseteq\mathcal T\), with an admissible kernel on each of
its intervals.  For each \(j\in\{0,1,2,3\}\), let
\(f_{j,s}:J\to\C\), \(0\le s\le d\), be measurable, and write
\(f_j=(f_{j,s})_{s=0}^d\).  Suppose that
\(\{0,\ldots,d\}\) is partitioned into \(W_j\ge1\) deterministic nonempty
consecutive blocks and that, on every block \(B\),
\[
 \mathcal V_B(f_j)(y)\le W_j^{-1/2}
 \quad\text{for almost every }y\in J.
\]
The four partitions may differ.  Extend the inputs by zero outside \(J\).
Then
\begin{equation}
 \sum_{I\in\mathcal I}|I|
  \bigl|H_I(f_{0,s(I)},f_{1,s(I)},f_{2,s(I)},f_{3,s(I)})\bigr|
 \le C_{\mathrm{loc}}(c_0,C_0)|J|.
 \label{eq:depth-local}
\end{equation}
The constant is independent of the lattice, root, tree, kernels, depth
range, blocks, and their numbers.
\end{theorem}

Sections~\ref{sec:rows}--\ref{sec:completion} prove
Theorem~\ref{thm:depth-local} from the specified analytic results of H.
Values are prescribed also on depths
without summation intervals; that full schedule is used in the proof.

The application below varies only slot zero.  The other three inputs use
the same functions at every depth and one block each, so their block condition
is simply the bound \(1\).  Theorem~\ref{thm:depth-local} applies uniformly
to translations of the lattice.

\subsection{Kernels and the continuous test}

For an averaging profile \(\varphi\) and \(a>0\),
Lemma~\ref{lem:transfer-kernel} constructs an admissible kernel whose
normalized spatial integral is \(\varphi_{2a}-\varphi_a\).
The unit-scale difference \(\varphi_2-\varphi\) has zero moments through
order three, hence a compactly supported fourth antiderivative.  The
construction multiplies a rescaling of this primitive by a smooth spatial
bump, then differentiates in the four cancellation directions.

\begin{lemma}[Exact realization of a smooth scale difference]
\label{lem:transfer-kernel}
For every averaging profile \(\varphi\) in
Definition~\ref{def:averaging-profile}, there are a dyadic integer
\(D_\varphi=2^d\), \(d\in\mathbb Z_{\geq0}\), a constant
\(C_{0,\varphi}\geq1\), and
\(\mathcal W_\varphi\in C_c^\infty(\mathbb R^2;\mathbb R)\) such that
the following holds.  For every \(a>0\) and every interval \(I\) of
length \(r=D_\varphi a\), with midpoint \(m_I\), the kernel
\[
 w_I(x,t)=\mathcal W_\varphi\left(\frac{x-m_I}{r},\frac{t}{r}\right)
\]
is admissible with \(c_0=1/8\) and \(C_0=C_{0,\varphi}\).  Moreover, for
every \(t\in\mathbb R\),
\begin{equation}
 \frac1r\int_{\mathbb R}\mathcal W_\varphi(z,t/r)\,dz
    =\varphi_{2a}(t)-\varphi_a(t).
 \label{eq:transfer-kernel-normalization}
\end{equation}
The parameters are fixed once \(\varphi\) is fixed.
\end{lemma}

\begin{proof}
Set \(b=\varphi_2-\varphi\).  For \(0\leq j\leq3\),
\[
 \int t^jb(t)\,dt=(2^j-1)\int t^j\varphi(t)\,dt=0,
\]
using the zero factor for \(j=0\) and \eqref{eq:profile-moments} for
the other moments.  Choose \(B\geq1\) with
\(\operatorname{supp}b\subseteq[-B,B]\), and define
\[
 g(t)=\frac16\int_{-\infty}^t(t-v)^3b(v)\,dv.
\]
It vanishes for \(t<-B\).  For \(t>B\), expansion of the cubic and the
four zero moments show that it again vanishes.  Differentiating gives
\(g^{(4)}=b\), so \(g\in C_c^\infty(\mathbb R)\) with support in
\([-B,B]\).  For \(0\leq j\leq3\),
\[
 g^{(j)}(t)=\frac1{(3-j)!}\int_{-\infty}^t(t-v)^{3-j}b(v)\,dv,
 \qquad
 \|g^{(j)}\|_\infty
       \leq\frac{(2B)^{4-j}}{(4-j)!}\|b\|_\infty.
\]
For \(n\geq4\), \(g^{(n)}=b^{(n-4)}\).  The derivative bounds for
\(\varphi\) therefore give bounds of the form
\eqref{eq:profile-derivatives} for all derivatives of \(g\), with a
constant depending on \(\varphi\).

Let \(\rho\) be the bump from Lemma~\ref{lem:ergodic-bump} and put
\(u(z)=8\rho(8z)\).  It has integral \(1\), support in \([-1/8,1/8]\),
and the same type of derivative bounds.  Choose a dyadic integer
\(D=D_\varphi>12B\).  Define
\[
 F(z,\sigma)=u(z)D^{-3}g(D\sigma),\qquad
 L_j=\partial_\sigma-j\partial_z,\qquad
 \mathcal W_\varphi=L_0L_1L_2L_3F.
\]
The operators commute, and differentiation does not enlarge support.
Thus, on the support of \(\mathcal W_\varphi\),
\[
 |z+j\sigma|\leq\frac18+\frac{3B}{D}<\frac38
       \qquad(0\leq j\leq3).
\]
For the scaled kernel this places every \(x+jt\) at distance strictly
greater than \(r/8\) from the endpoints of \(I\).

For \(j=0,1,2,3\), put \(G_j=\prod_{\ell\neq j}L_\ell F\).  The chain
rule gives, for every \(v,t\in\mathbb R\),
\[
 w_I(v-jt,t)
   =r\,\frac{d}{dt}
       G_j\left(\frac{v-jt-m_I}{r},\frac{t}{r}\right).
\]
The differentiated function is compactly supported in \(t\), so the
integral of its derivative is zero.  This proves each of the four
cancellation conditions for every \(v\).

Expansion of the operators gives
\[
 \mathcal W_\varphi(z,\sigma)
   =D u(z)g^{(4)}(D\sigma)-6u'(z)g^{(3)}(D\sigma)
      +11D^{-1}u''(z)g''(D\sigma)-6D^{-2}u'''(z)g'(D\sigma).
\]
A derivative \(\partial_z^p\partial_\sigma^q\) of total order \(n\)
is a sum of four products of derivatives of \(u,g\) of orders at most
\(n+4\), with factors bounded by a constant times \(D^{n+1}\).
For fixed \(D\), their derivative bounds yield
\[
 \|\partial_z^p\partial_\sigma^q\mathcal W_\varphi\|_\infty
   \leq\bigl(C_{0,\varphi}(n+2)\bigr)^{C_{0,\varphi}(n+2)}
\]
for one sufficiently large \(C_{0,\varphi}\) and all \(p,q\).
Rescaling gives the required factor \(r^{-n}\) for derivatives in
\(x,t\).  Finally, integration in \(z\) kills all three terms containing
derivatives of \(u\), and hence
\[
 \int_{\mathbb R}\mathcal W_\varphi(z,\sigma)\,dz
       =Dg^{(4)}(D\sigma)=Db(D\sigma).
\]
As \(r=Da\), division by \(r\) gives
\(a^{-1}b(t/a)=\varphi_{2a}(t)-\varphi_a(t)\).
\end{proof}

For the rest of this section fix an averaging profile \(\varphi\), and
write \(D=D_\varphi\), \(b=\varphi_2-\varphi\), and
\[
 \Delta_k(t)=\varphi_{2^{k+1}}(t)-\varphi_{2^k}(t)
            =2^{-k}b(2^{-k}t),\qquad k\in\mathbb Z.
\]
Constants with subscript \(\varphi\) may depend on this profile, the
fixed choices in Lemma~\ref{lem:transfer-kernel}, and
\(C_{\mathrm{loc}}(1/8,C_{0,\varphi})\).  They are independent of
the scales, blocks, inputs, and probability system.

\begin{lemma}[A continuous test obtained by translation averaging]
\label{lem:transfer-continuous}
For the fixed profile \(\varphi\) and \(D=D_\varphi\), there is
\(C_\varphi^{\mathrm c}<\infty\) with the following property.
Let \(k_-\leq k_+\) be integers, and partition
\(\{k_-,\ldots,k_+\}\) into \(w\geq1\) deterministic nonempty consecutive
blocks.  Let \(K\subset\mathbb R\) be an interval of length \(L>0\).
Let \(f_1,f_2,f_3:\mathbb R\to\mathbb C\) and
\(f_{0,k}:\mathbb R\to\mathbb C\), \(k_-\leq k\leq k_+\), be measurable
and vanish almost everywhere outside \(K\).  Suppose
\(|f_j|\leq1\) almost everywhere for \(j=1,2,3\), and
\(\mathcal V_B(f_0)(x)\leq w^{-1/2}\) for almost every \(x\in\mathbb R\)
and every block \(B\).
Then
\begin{equation}
 \left|\sum_{k=k_-}^{k_+}\int_{\mathbb R^2}
       f_{0,k}(x)\prod_{j=1}^3f_j(x+jt)\Delta_k(t)\,dx\,dt\right|
       \leq C_\varphi^{\mathrm c}(L+D2^{k_+}).
 \label{eq:transfer-continuous}
\end{equation}
\end{lemma}

\begin{proof}
Set \(r_k=D2^k\) and \(R=r_{k_+}\).  For \(0\leq h<R\), use the one
translated dyadic lattice
\[
 \mathcal D_h
   =\{[h+m2^q,h+(m+1)2^q):m,q\in\mathbb Z\}.
\]
Every \(r_k\) is a length in this lattice and divides \(R\).
Use the kernels of Lemma~\ref{lem:transfer-kernel}, and let \(H_I^h\)
denote the local form with inputs \(f_{0,k},f_1,f_2,f_3\) when
\(|I|=r_k\).  The zeroth coordinate \(x\) of a supported kernel lies
in \(I\).  Hence every nonzero summand belongs to a root interval of
length \(R\) meeting \(K\).  The total length of these finitely many
roots is at most \(L+2R\).

In each root, the intervals at the specified lengths lie in a finite
tree of depth \(k_+-k_-\).  Restrict the inputs to that root; all forms
are unchanged because every progression coordinate on the support lies
in the summation interval.  The fixed slots use constant depth
schedules with one block.  For slot zero, depth is \(s=k_+-k\);
reversing the scale order reverses the block order and the order inside
each block without changing \(\mathcal V_B\).  All hypotheses of
Theorem~\ref{thm:depth-local} therefore hold on this root.  Summing its
conclusion over the roots gives
\begin{equation}
 \left|\sum_{k=k_-}^{k_+}
       \sum_{\substack{I\in\mathcal D_h\\|I|=r_k}}r_kH_I^h\right|
 \leq\sum_{k=k_-}^{k_+}
       \sum_{\substack{I\in\mathcal D_h\\|I|=r_k}}r_k|H_I^h|
 \leq C_{\mathrm{loc}}(1/8,C_{0,\varphi})(L+2R).
 \label{eq:transfer-before-averaging}
\end{equation}

For \(r=r_k\), the spatial sum below is \(r\)-periodic in \(h\).
Since \(r\) divides \(R\), averaging and unfolding the translates gives,
for every \(x,t\in\mathbb R\),
\begin{align*}
 \frac1R\int_0^R
    \sum_{\substack{I\in\mathcal D_h\\|I|=r}}r^{-1}w_I(x,t)\,dh
 &=\frac1r\int_0^r\sum_{m\in\mathbb Z}r^{-1}
      \mathcal W_\varphi\left(
        \frac{x-h-(m+1/2)r}{r},\frac{t}{r}\right)\,dh\\
 &=\frac1r\int_{\mathbb R}\mathcal W_\varphi(z,t/r)\,dz
   =\Delta_k(t).
\end{align*}
Here the last equality is \eqref{eq:transfer-kernel-normalization}.
The normalization of the form is
\[
 rH_I^h=r^{-1}\int_{\mathbb R^2}
      w_I(x,t)f_{0,k}(x)\prod_{j=1}^3f_j(x+jt)\,dx\,dt.
\]
It follows that averaging the first sum in
\eqref{eq:transfer-before-averaging} over \(h\in[0,R)\) gives exactly
the sum of integrals in \eqref{eq:transfer-continuous}.  The finite
scale range, bounded common support of the inputs, and kernel support
justify the exchanges of sums and integrals.  Taking the absolute value
after averaging proves \eqref{eq:transfer-continuous}, with
\(C_\varphi^{\mathrm c}=2C_{\mathrm{loc}}(1/8,C_{0,\varphi})\) as one
possible choice.
\end{proof}

\subsection{Planting integer sequences and transferring to orbits}

For sequences on \(\mathbb Z\), support in an integer interval means
vanishing at every integer outside it.  We measure such an interval by
its cardinality.

\begin{lemma}[The corresponding test on integer sequences]
\label{lem:transfer-discrete}
For the fixed profile \(\varphi\) and \(D=D_\varphi\), there is
\(C_\varphi^{\mathrm d}<\infty\) with the following property.
Let \(0\leq k_-\leq k_+\) be integers, and partition
\(\{k_-,\ldots,k_+\}\) into \(w\geq1\) deterministic nonempty consecutive
blocks.  Let \(\mathcal K=\{q,\ldots,q+L-1\}\), with \(q\in\mathbb Z\)
and integer \(L\geq1\).  Suppose
\(f_1,f_2,f_3:\mathbb Z\to\mathbb C\) and
\(f_{0,k}:\mathbb Z\to\mathbb C\), \(k_-\leq k\leq k_+\), are supported
in \(\mathcal K\), with \(|f_j(m)|\leq1\) for every \(m\) and
\(j=1,2,3\).  Suppose also that
\(\mathcal V_B(f_0)(m)\leq w^{-1/2}\) for every \(m\in\mathbb Z\) and
every block.  Then
\begin{equation}
 \left|\sum_{k=k_-}^{k_+}\sum_{m,n\in\mathbb Z}
    f_{0,k}(m)\prod_{j=1}^3f_j(m+jn)\Delta_k(n)\right|
       \leq C_\varphi^{\mathrm d}(L+D2^{k_+}).
 \label{eq:transfer-discrete}
\end{equation}
\end{lemma}

\begin{proof}
Fix \(e=1/100\).  For any one of the input sequences, define
\[
 \widetilde f(y)=\sum_{m\in\mathbb Z}
       f(m)\mathbf 1_{[m-e,m+e)}(y).
\]
These small intervals are disjoint.  The planted inputs are supported
in \([q-e,q+L-1+e]\), which is contained in a real interval of length
\(L\).  Their size and block conditions persist pointwise, because
every \(y\) belongs to at most one of the half-open planted intervals.

If the product of the four planted inputs is nonzero at \((x,t)\),
let \(q_j\) be its unique integer centers with \(|x+jt-q_j|\leq e\).
We discard the null set of boundary equalities.  Put
\[
 m=q_0,\qquad n=q_1-q_0,\qquad \xi=x-m,\qquad\eta=t-n.
\]
The first two coordinates give \(|\xi|\leq e\) and \(|\eta|\leq2e\).
For \(j=2,3\),
\[
 |q_j-(m+jn)|
 \leq |q_j-(x+jt)|+|\xi+j\eta|
 \leq(2j+2)e<1.
\]
This integer difference is zero.  Conversely, the coordinates lie in
the four small intervals centered at \(m+jn\) precisely on the
translate of
\[
 Q_e=\{(\xi,\eta):|\xi+j\eta|\leq e,\ 0\leq j\leq3\},
\]
up to boundaries.  The two middle inequalities follow from the endpoint
inequalities by convexity.  The map
\((\xi,\eta)\mapsto(\xi,\xi+3\eta)\) has determinant \(3\), so
\[
 |Q_e|=\frac{4e^2}{3}=:A_e>0.
\]
Distinct integer pairs give disjoint regions up to null boundaries,
because their zeroth and first centers are unique.  Therefore the
continuous planted integral at scale \(k\) is exactly
\[
 I_k=\sum_{m,n\in\mathbb Z}
       f_{0,k}(m)\prod_{j=1}^3f_j(m+jn)
       \int_{Q_e}\Delta_k(n+\eta)\,d\xi\,d\eta.
\]

Choose \(B_b\geq1\) with \(\operatorname{supp}b\subseteq[-B_b,B_b]\).
At scale \(a=2^k\geq1\), the mean value inequality gives
\[
 |\Delta_k(n+\eta)-\Delta_k(n)|
       \leq2e\|b'\|_\infty a^{-2}\qquad((\xi,\eta)\in Q_e).
\]
If the difference is nonzero, then \(|n|\leq B_ba+2e\), so there are at
most \((2B_b+4e+1)a\) possible integers \(n\).  There are at most \(L\)
contributing integers \(m\), and every coefficient product has absolute
value at most \(1\), since
\(\mathcal V_B(f_0)(m)\leq w^{-1/2}\leq1\).  Consequently, with
\[
 D_k=\sum_{m,n\in\mathbb Z}
       f_{0,k}(m)\prod_{j=1}^3f_j(m+jn)\Delta_k(n),
\]
we have
\[
 |I_k-A_eD_k|
 \leq A_e\,2e\|b'\|_\infty(2B_b+4e+1)L\,2^{-k}.
\]
The restriction \(k\geq0\) makes the errors summable, since
\(\sum_{k=k_-}^{k_+}2^{-k}\leq2\).  Apply
Lemma~\ref{lem:transfer-continuous} to the planted inputs, sum this error
bound, and divide by \(A_e\).  The resulting error is at most a constant
depending only on \(\varphi\) times \(L\), proving
\eqref{eq:transfer-discrete}.
\end{proof}

\begin{proof}[Proof of Proposition~\ref{prop:orbit-oscillation}]
We first obtain the corresponding finite estimate on integer sequences.
Let \(\mathcal K\) be an integer interval of cardinality \(L\geq1\),
let \(\mathcal E\subseteq\mathcal K\) be a nonempty integer interval, and
let \(f_1,f_2,f_3:\mathbb Z\to\mathbb C\) be supported in \(\mathcal K\)
and bounded in absolute value by \(1\).  Put
\[
 \mathsf V_k(m)=\sum_{n\in\mathbb Z}\varphi_{2^k}(n)
                    \prod_{j=1}^3f_j(m+jn),
       \qquad k\geq0,\quad m\in\mathbb Z.
\]
Fix deterministic integers \(0\leq k_1<\cdots<k_{w+1}\), with \(w\geq1\).
For \(m\in\mathcal E\), let \(\ell_i(m)\) be the smallest index attaining
the maximum of \(|\mathsf V_\ell(m)-\mathsf V_{k_i}(m)|\) over
\(k_i\leq\ell\leq k_{i+1}\).  Set
\[
 d_i(m)=\mathsf V_{\ell_i(m)}(m)-\mathsf V_{k_i}(m),\qquad
 \alpha_i(m)=
 \begin{cases}\overline{d_i(m)}/|d_i(m)|,&d_i(m)\neq0,\\
 0,&d_i(m)=0.\end{cases}
\]
Thus \(\alpha_i(m)d_i(m)=|d_i(m)|\), also for complex inputs.  In the
block \(k_i\leq k<k_{i+1}\), define
\[
 f_{0,k}(m)=
 \begin{cases}
  \alpha_i(m)/(4\sqrt w),&m\in\mathcal E\text{ and }k<\ell_i(m),\\
  0,&\text{otherwise}.
 \end{cases}
\]
On that block the supremum is at most \((4\sqrt w)^{-1}\), and there
is at most one jump of that size.  Hence
\(\mathcal V_B(f_0)(m)\leq(2\sqrt w)^{-1}\leq w^{-1/2}\).
The block endpoints are deterministic; only the prefix lengths and
phases depend on \(m\).

The scale differences telescope exactly:
\begin{align*}
 &\sum_{k=k_1}^{k_{w+1}-1}\sum_{m,n\in\mathbb Z}
       f_{0,k}(m)\prod_{j=1}^3f_j(m+jn)\Delta_k(n)\\
 &=\frac1{4\sqrt w}\sum_{m\in\mathcal E}\sum_{i=1}^w
       \alpha_i(m)\bigl(\mathsf V_{\ell_i(m)}(m)
                              -\mathsf V_{k_i}(m)\bigr)\\
 &=\frac1{4\sqrt w}\sum_{m\in\mathcal E}\sum_{i=1}^w
       \max_{k_i\leq\ell\leq k_{i+1}}
            |\mathsf V_\ell(m)-\mathsf V_{k_i}(m)|.
\end{align*}
The final expression is nonnegative real.  Lemma~\ref{lem:transfer-discrete}
applies to the blocks just constructed.  Its largest difference scale
is \(2^{k_{w+1}-1}\leq2^{k_{w+1}}\).  We obtain
\begin{equation}
 \sum_{m\in\mathcal E}\sum_{i=1}^w
       \max_{k_i\leq\ell\leq k_{i+1}}
            |\mathsf V_\ell(m)-\mathsf V_{k_i}(m)|
 \leq 4C_\varphi^{\mathrm d}\sqrt w\,
                  (L+D2^{k_{w+1}}).
 \label{eq:transfer-sequence-oscillation}
\end{equation}

We now use the finite-orbit form of Calder\'on's transference principle;
compare the explicit cutoff arguments in \cite[Section~2B]{Bourgain1989}
and \cite{KMT2022}, as well as \cite{Calderon1968}.
Let the system and functions be as in
Proposition~\ref{prop:orbit-oscillation}, and retain the fixed endpoint
list.  Choose \(B_\varphi\geq1\) with
\(\operatorname{supp}\varphi\subseteq[-B_\varphi,B_\varphi]\), and set
\(Q=\lceil3B_\varphi2^{k_{w+1}}\rceil\).  For each integer \(M\geq1\)
and \(x\in X\), use the finite sequences
\[
 g_{j,x}(r)=
 \begin{cases}
  f_j(T^r x),&1-Q\leq r\leq M+Q,\\
  0,&\text{otherwise},
 \end{cases}
 \qquad r\in\mathbb Z,\quad j=1,2,3.
\]
Their common support interval has cardinality \(M+2Q\), and contains
the testing interval \(\{1,\ldots,M\}\).  If \(1\leq m\leq M\),
\(k_1\leq\ell\leq k_{w+1}\), and \(\varphi_{2^\ell}(n)\neq0\), then
\(|jn|\leq3B_\varphi2^{k_{w+1}}\leq Q\).  Thus the sequence average
is exactly
\[
 \sum_n\varphi_{2^\ell}(n)\prod_{j=1}^3g_{j,x}(m+jn)
    =V_\ell^\varphi(T^m x).
\]
Apply \eqref{eq:transfer-sequence-oscillation} to these sequences for each
\(x\).  The resulting finite maxima are measurable and bounded, as observed
after \eqref{eq:smooth-uniform-bound}.  Integrating and using invariance of
each \(T^m\) gives
\[
 M\sum_{i=1}^w\int_X
    \max_{k_i\leq\ell\leq k_{i+1}}
        |V_\ell^\varphi-V_{k_i}^\varphi|\,d\mu
 \leq 4C_\varphi^{\mathrm d}\sqrt w
        (M+2Q+D2^{k_{w+1}}).
\]
Divide by \(M\) and let \(M\to\infty\), keeping the endpoint
list, hence \(Q\), fixed.  We obtain \eqref{eq:orbit-oscillation} with
\(C_\varphi=4C_\varphi^{\mathrm d}\).  This finite transfer uses only
measurability and measure preservation of the integer powers of \(T\).
\end{proof}
